\documentclass[12pt, a4paper]{article}

% --- PACCHETTI ---
\usepackage[utf8]{inputenc}
\usepackage[italian]{babel}
\usepackage{amsmath, amsfonts, amssymb}
\usepackage{graphicx}
\usepackage{hyperref}
\usepackage{booktabs}
\usepackage{float}
\usepackage{caption}
\usepackage{subcaption}
\usepackage{geometry}
\usepackage{setspace}
\usepackage[protrusion=true, expansion=true, patches]{microtype}
\usepackage{xcolor}
\usepackage{enumitem}
\usepackage{array}
\usepackage{parskip}
\usepackage{siunitx}
\usepackage{tcolorbox}
\tcbuselibrary{skins, breakable}

\geometry{margin=2.5cm}
\onehalfspacing

\hypersetup{
  colorlinks = true,
  linkcolor  = blue,
  citecolor  = blue,
  urlcolor   = blue,
}

% --- Comandi utili ---
\newcommand{\code}[1]{\texttt{#1}}
\newcommand{\pkg}[1]{\textsf{#1}}

% --- Box definizione ---
\newtcolorbox{defbox}[1][]{
  enhanced, breakable,
  colback  = blue!5!white,
  colframe = blue!40!black,
  fonttitle = \bfseries,
  title = {#1},
  left=6pt, right=6pt, top=4pt, bottom=4pt
}

% --- TITOLO ---
\title{
  \vspace{-1cm}
  {\Large Università degli Studi di Padova} \\[0.4cm]
  {\large Corso di Metodi Statistici per Big Data} \\[1cm]
  \rule{\linewidth}{0.5pt} \\[0.4cm]
  {\LARGE \textbf{Analisi Statistica del Mercato Cinematografico}} \\[0.3cm]
  {\large Text mining e modelli predittivi sulle sinossi TMDB} \\
  \rule{\linewidth}{0.5pt}
}
\author{
  Andrea Fior \texttt{2111172} \and
  Massimo Dovidio \texttt{2117366} \and
  Matteo Mauggeri \texttt{2103500}
}
\date{Anno Accademico 2025/2026}

% ================================================================
\begin{document}
\maketitle
\thispagestyle{empty}
\newpage
\tableofcontents
\newpage

% ================================================================
\section{Introduzione}

Il presente studio analizza le dinamiche del mercato cinematografico
attraverso un approccio che integra metodi statistici con tecniche
di \emph{text mining}.
Partendo da un dataset di \textbf{10\,000 film} estratto tramite le
API di The Movie Database (TMDB), l'analisi si propone di
investigare in quale misura il \emph{contenuto narrativo} di un
film --- sintetizzato dalla sua sinossi --- sia in grado di
spiegare e prevedere il gradimento del pubblico, misurato tramite
il punteggio medio (\code{vote\_average}) e la variabile binaria
derivata \code{high\_rated}.

% ----------------------------------------------------------------
\subsection{Descrizione del Dataset}

Il dataset è stato ottenuto tramite le API pubbliche di
\href{https://www.themoviedb.org/}{The Movie Database (TMDB)}.
Dopo aver rimosso le osservazioni con meno di 10 voti, prive di
sinossi o con data di uscita successiva al 31~dicembre~2025,
il dataset finale conta \textbf{4\,867 film}.

\begin{table}[H]
\centering
\caption{Variabili del dataset TMDB utilizzate nell'analisi}
\smallskip
\begin{tabular}{@{}llp{7cm}@{}}
\toprule
\textbf{Variabile} & \textbf{Tipo} & \textbf{Descrizione} \\
\midrule
\code{title}              & Testuale    & Titolo della pellicola \\
\code{overview}           & Testuale    & Sinossi della trama (input principale) \\
\code{original\_language} & Categoriale & Lingua originale \\
\code{genre\_ids}         & Categoriale & Identificatori numerici dei generi \\
\code{vote\_average}      & Numerica    & Punteggio medio utenti TMDB (0--10) \\
\code{vote\_count}        & Numerica    & Numero totale di voti ricevuti \\
\code{popularity}         & Numerica    & Indice di popolarità TMDB \\
\code{release\_date}      & Data        & Data di uscita \\
\bottomrule
\end{tabular}
\end{table}

Sono state create le seguenti \textbf{variabili derivate}:

\begin{itemize}[leftmargin=1.8em]
  \item \code{genre}: genere principale decodificato da
        \code{genre\_ids} tramite la tabella di lookup TMDB
        (18~generi più la categoria \textit{Other});
  \item \code{high\_rated}: variabile binaria con livelli
        \textit{Alto} (\code{vote\_average} $\geq 7$,
        circa il \textbf{22.8\%} dei film) e \textit{Basso}
        (restante 77.2\%);
  \item \code{year}: anno di uscita estratto da
        \code{release\_date};
  \item \code{n\_parole}, \code{n\_caratteri}: lunghezza
        della sinossi in parole e caratteri, calcolate
        \emph{prima} del preprocessing.
\end{itemize}

% ----------------------------------------------------------------
\subsection{Percorso Metodologico}

Il progetto segue la pipeline seguente:

\begin{enumerate}[leftmargin=1.8em]
  \item \textbf{Preprocessing} del testo delle sinossi.
  \item \textbf{Analisi esplorativa}: frequenze, wordcloud,
        bigrammi, lunghezza per genere, distribuzione del voto.
  \item \textbf{Topic Modeling (LDA)} per la struttura tematica
        latente delle sinossi.
  \item \textbf{Topic Modeling Supervisionato (sLDA)} per
        collegare i topic al voto.
  \item \textbf{Regressione penalizzata} (Lasso, Ridge,
        Elastic Net) per prevedere \code{vote\_average}.
  \item \textbf{Classificazione} di \code{high\_rated} tramite
        LASSO Logistico e Random Forest.
\end{enumerate}


% ================================================================
\section{Preprocessing del Testo}

Le sinossi vengono sottoposte a una pipeline di pulizia
implementata con il pacchetto \pkg{tm} e funzioni personalizzate:

\begin{enumerate}[leftmargin=1.8em]
  \item Rimozione di URL, hashtag, caratteri di controllo e
        spazi ridondanti.
  \item Conversione in minuscolo.
  \item Rimozione delle stopword inglesi (lista \texttt{SMART},
        applicata due volte per coprire residui
        post-punteggiatura).
  \item Rimozione della punteggiatura e dei numeri.
  \item Rimozione di un dizionario personalizzato di termini
        non informativi per il dominio cinematografico:
        \textit{film, movie, story, life, one, finds, can, will,
        two, new, world, set, takes, make, time, back,
        love, family}.
  \item Filtraggio dei token di lunghezza $\leq 2$ caratteri.
\end{enumerate}

Il risultato è il dataframe \code{tidy\_overview}, che mantiene
le colonne \code{film\_id} e \code{genre} per tutte le analisi
successive.
Il vocabolario del corpus presenta il \textbf{38.1\%} di
\textit{hapax} (parole con frequenza~1) e il \textbf{78.4\%}
di parole con frequenza $\leq 5$, evidenziando una distribuzione
lessicale fortemente asimmetrica tipica dei testi brevi.


% ================================================================
\section{Analisi Esplorativa del Testo}

% ----------------------------------------------------------------
\subsection{Distribuzione del Voto}

La Figura~\ref{fig:distrib_voti} mostra la distribuzione di
\code{vote\_average} sull'intero dataset, suddivisa nelle due
classi \textit{Alto} ($\geq 7$) e \textit{Basso} ($< 7$).
La distribuzione risulta asimmetrica negativa, con il picco
intorno a 6.5. Solo il \textbf{22.8\%} dei film supera la soglia
di 7, evidenziando lo sbilanciamento strutturale che condizionerà
le scelte modellistiche nelle sezioni successive.

\begin{figure}[H]
  \centering
  \includegraphics[width=0.85\textwidth]{votesdistr.png}
  \caption{Distribuzione di \code{vote\_average} suddivisa per classe.
           In blu i film con voto Alto ($\geq 7$), in rosso quelli Basso.}
  \label{fig:distrib_voti}
\end{figure}

% ----------------------------------------------------------------
\subsection{Parole più Frequenti}

Le 20 parole più frequenti sono visualizzate con un grafico a
barre orizzontali con intensità cromatica proporzionale alla
frequenza (Figura~\ref{fig:top20}); una wordcloud ne riassume
le prime 100.

\begin{figure}[H]
  \centering
  \begin{subfigure}[t]{0.80\textwidth}
    \includegraphics[width=\linewidth]{20parole.png}
    \caption{Top 20 parole per frequenza}
  \end{subfigure}
  \hfill
  \begin{subfigure}[t]{0.80\textwidth}
    \includegraphics[width=\linewidth]{wordcloud1.png}
    \caption{Wordcloud delle 100 parole più frequenti}
  \end{subfigure}
  \caption{Frequenza delle parole nelle sinossi dopo preprocessing}
  \label{fig:top20}
\end{figure}

Tra le parole più ricorrenti emergono termini narrativi
universali (\textit{man, woman, father, mother, young, death})
e vocaboli d'azione (\textit{police, murder, escape, mission}),
coerentemente con la composizione per genere del dataset.

% ----------------------------------------------------------------
\subsection{Parole Distintive per Genere (Barplot)}

Per i sei generi principali (\textit{Action, Drama, Comedy,
Horror, Thriller, Romance}) viene riportata la top-10 delle
parole più frequenti in ciascun genere, visualizzate con grafici
a barre sfaccettati (Figura~\ref{fig:barplot_generi}).
Questo approccio sostituisce la wordcloud per genere e consente
un confronto quantitativo diretto delle frequenze.

\begin{figure}[H]
  \centering
  \includegraphics[width=\textwidth]{barplot_generi.pdf}
  \caption{Top 10 parole per frequenza nei sei generi principali.
           Ogni pannello è ordinato in modo decrescente per frequenza.}
  \label{fig:barplot_generi}
\end{figure}

I pattern lessicali sono nettamente distinti: il genere
\textit{Horror} è dominato da termini come \textit{haunted,
evil, night}; il \textit{Drama} da \textit{mother, father,
relationship}; l'\textit{Action} da \textit{mission, police,
escape}. La separazione lessicale conferma che la variabile
\code{genre} cattura informazione testuale reale.

% ----------------------------------------------------------------
\subsection{Bigrammi}

L'analisi dei bigrammi considera coppie di parole consecutive,
filtrate per stopword e dizionario personalizzato.
Il grafico a barre (Figura~\ref{fig:bigrammi}) mostra i
20~bigrammi più frequenti.

\begin{figure}[H]
  \centering
    \includegraphics[width=\linewidth]{bigrams.png}
    \caption{Top 20 bigrammi per frequenza}
    \label{fig:bigrammi}
\end{figure}

I bigrammi rivelano costruzioni tipiche del linguaggio
cinematografico --- come \textit{serial killer},
\textit{young man}, \textit{high school}, \textit{special
agent} --- che il conteggio di singoli token non può catturare.

% ----------------------------------------------------------------
\subsection{Distribuzione del Voto per Genere}

La Figura~\ref{fig:voto_genere} mostra la distribuzione di
\code{vote\_average} per ciascun genere tramite boxplot.
La linea tratteggiata a 7 indica la soglia che definisce la
classe \textit{Alto}. I generi \textit{Animation},
\textit{War} e \textit{Documentary} presentano mediane più alte,
mentre \textit{Horror}, \textit{Mystery} e \textit{Thriller}
mostrano distribuzioni più basse e variabili.

\begin{figure}[H]
  \centering
  \includegraphics[width=\textwidth]{votixgenere.png}
  \caption{Distribuzione di \code{vote\_average} per genere.
           La linea tratteggiata rossa indica la soglia $\geq 7$.}
  \label{fig:voto_genere}
\end{figure}

% ----------------------------------------------------------------
\subsection{Lunghezza della Trama e Voto}

La Figura~\ref{fig:lun_voto} esplora la relazione tra numero
di parole nella sinossi e voto medio, colorando i film per
classe di gradimento.

\begin{figure}[H]
  \centering
  \begin{subfigure}[t]{0.48\textwidth}
    \includegraphics[width=\linewidth]{plot_lunghezza_voto_scatter.pdf}
    \caption{Scatterplot con curva GAM}
  \end{subfigure}
  \hfill
  \begin{subfigure}[t]{0.48\textwidth}
    \includegraphics[width=\linewidth]{densities.png}
    \caption{Densità per classe di voto}
  \end{subfigure}
  \caption{Relazione tra lunghezza della sinossi e gradimento}
  \label{fig:lun_voto}
\end{figure}


% ================================================================
\section{Topic Modeling Non Supervisionato (LDA)}

\subsection{Definizione e Metodo}

\begin{defbox}[Latent Dirichlet Allocation (LDA)]
La \textbf{Latent Dirichlet Allocation} (Blei, Ng e Jordan, 2003)
è un modello generativo bayesiano per collezioni di documenti.
Ogni documento viene rappresentato come una \emph{miscela} di
$K$~topic latenti, e ogni topic come una distribuzione di
probabilità sul vocabolario. I parametri del modello vengono
stimati tramite inferenza bayesiana approssimata.

Formalmente, per ciascun documento $d$ e ciascuna parola $w_i$:
\[
  p(w_i \mid d) = \sum_{k=1}^{K} p(w_i \mid z_i = k)\, p(z_i = k \mid d)
\]
dove $\theta_d \sim \mathrm{Dir}(\alpha)$ è la distribuzione dei
topic nel documento e $\phi_k \sim \mathrm{Dir}(\eta)$ è la
distribuzione delle parole nel topic $k$.
La stima è condotta tramite \textbf{Gibbs Sampling collassato}
(pacchetto \pkg{lda}), che campiona iterativamente l'assegnazione
di topic per ogni token.

L'LDA è \emph{non supervisionato}: i topic sono appresi
esclusivamente dalla struttura co-occorrenziale del testo,
senza usare alcuna variabile risposta.
\end{defbox}

\begin{table}[H]
\centering
\caption{Iperparametri del modello LDA}
\smallskip
\begin{tabular}{@{}lll@{}}
\toprule
\textbf{Parametro} & \textbf{Valore} & \textbf{Significato} \\
\midrule
$K$          & 6   & Numero di topic \\
$\alpha$     & 0.1 & Prior di Dirichlet sui topic (sparse) \\
$\eta$       & 0.1 & Prior di Dirichlet sulle parole (sparse) \\
Iterazioni   & 500 & Iterazioni totali del sampler \\
Burnin       & 100 & Iterazioni di riscaldamento scartate \\
\bottomrule
\end{tabular}
\end{table}

Il vocabolario è ristretto alle parole presenti in almeno
5~sinossi distinte; nessun film risulta con documento vuoto
dopo il filtraggio (4\,867 documenti usati).

\subsection{Topic Identificati}

\begin{table}[H]
\centering
\caption{Top 8 parole per topic (LDA, $K=6$, \texttt{seed = 12345})}
\smallskip
\setlength{\tabcolsep}{5pt}
\begin{tabular}{@{}cllllll@{}}
\toprule
\textbf{\#} &
\textbf{T1: Sci-Fi \&} &
\textbf{T2: Romance \&} &
\textbf{T3: Crimine \&} &
\textbf{T4: Dramma} &
\textbf{T5: Thriller \&} &
\textbf{T6: Guerra \&} \\
& \textbf{Avventura} &
\textbf{Scuola} &
\textbf{Città} &
\textbf{Familiare} &
\textbf{Mistero} &
\textbf{Azione} \\
\midrule
1 & mother  & police   & school   & mother  & war       & earth   \\
2 & woman   & agent    & meets    & woman   & american  & planet  \\
3 & relationship & gang & friend  & haunted & british   & mission \\
4 & child   & prison   & day      & night   & german    & evil    \\
5 & day     & murder   & beautiful & mysterious & village & human  \\
6 & sexual  & drug     & night    & past    & black     & alien   \\
7 & lives   & criminal & plan     & couple  & west      & team    \\
8 & home    & crime    & play     & home    & soldiers  & quest   \\
\bottomrule
\end{tabular}
\end{table}

La distribuzione dei film per topic risulta:
\textit{Dramma Familiare}~888,
\textit{Crimine \& Città}~865,
\textit{Romance \& Scuola}~861,
\textit{Sci-Fi \& Avventura}~841,
\textit{Guerra \& Azione}~746,
\textit{Thriller \& Mistero}~666.

\subsection{Risultati Grafici}

\begin{figure}[H]
  \centering
  \includegraphics[width=0.85\textwidth]{plot_lda_convergenza.pdf}
  \caption{Convergenza del Gibbs Sampler: log-likelihood per
           iterazione. La linea tratteggiata rossa indica la
           fine del burnin (100~iterazioni).}
  \label{fig:lda_conv}
\end{figure}

\begin{figure}[H]
  \centering
  \includegraphics[width=\textwidth]{plot_lda_topic_genere.pdf}
  \caption{Distribuzione percentuale dei topic dominanti
           per i sei generi principali}
  \label{fig:lda_genere}
\end{figure}

\begin{figure}[H]
  \centering
  \begin{subfigure}[t]{0.48\textwidth}
    \includegraphics[width=\linewidth]{plot_lda_voto_topic.pdf}
    \caption{Distribuzione di \code{vote\_average} per topic}
  \end{subfigure}
  \hfill
  \begin{subfigure}[t]{0.48\textwidth}
    \includegraphics[width=\linewidth]{plot_lda_pop_voto.pdf}
    \caption{Popolarità media vs voto medio per topic}
  \end{subfigure}
  \caption{Relazione tra topic dominante, gradimento e popolarità}
  \label{fig:lda_voto}
\end{figure}

Il grafico di dispersione (Figura~\ref{fig:lda_voto}b) mostra
che i topic legati ad azione e fantascienza tendono ad avere
popolarità media più elevata, mentre quelli drammatici
presentano voti mediamente superiori, suggerendo una tensione
tra apprezzamento critico e successo commerciale.


% ================================================================
\section{Topic Modeling Supervisionato (sLDA)}

\subsection{Definizione, Differenze e Limiti rispetto a LDA}

\begin{defbox}[Supervised Latent Dirichlet Allocation (sLDA)]
La \textbf{Supervised LDA} (Blei e McAuliffe, 2007) estende
l'LDA introducendo una variabile risposta $y_d$ per ciascun
documento. In questo studio $y_d = \code{vote\_average}_d$.
Il modello assume che la variabile risposta dipenda linearmente
dalla proporzione di topic del documento:
\[
  y_d \mid \bar{\theta}_d, \eta, \sigma^2
  \;\sim\;
  \mathcal{N}\!\left(\eta^\top \bar{\theta}_d,\; \sigma^2\right)
\]
dove $\bar{\theta}_d$ è il vettore medio delle assegnazioni di
topic nel documento $d$ e $\eta \in \mathbb{R}^K$ è il vettore
dei coefficienti di regressione.

La stima avviene tramite un algoritmo \textbf{EM variazionale}
che alterna:
\begin{itemize}[leftmargin=1.8em, topsep=2pt]
  \item \textbf{E-step}: aggiorna i parametri variazionali
        $(\gamma_d, \phi_{d,n})$ tenendo fissi $\eta$ e $\sigma^2$;
  \item \textbf{M-step}: aggiorna $\eta$ e $\sigma^2$ via
        regressione lineare sui $\bar{\theta}_d$ stimati.
\end{itemize}
\end{defbox}

\paragraph{Differenze rispetto all'LDA non supervisionato.}
La differenza fondamentale è che in LDA i topic vengono appresi
\emph{solo} dalla distribuzione delle parole nel corpus, mentre
in sLDA i topic vengono appresi \emph{congiuntamente} dalla
struttura del testo e dal segnale della variabile risposta.
In pratica, sLDA tende a separare i topic in modo più utile
per la predizione: topic che si rivelano irrilevanti per $y$
possono essere compressi, e vengono preferiti topic che
``spiegano'' la varianza del voto.
Questo porta a rappresentazioni $\theta_d$ più informative ai
fini regressivi, a costo di una minore interpretabilità tematica
\emph{pura} dei topic.

Nella Tabella~\ref{tab:lda_vs_slda} si riassumono le principali
differenze tra i due approcci.

\begin{table}[H]
\centering
\caption{Confronto tra LDA e sLDA}
\label{tab:lda_vs_slda}
\smallskip
\begin{tabular}{@{}p{4.2cm}p{4.5cm}p{4.5cm}@{}}
\toprule
\textbf{Caratteristica} & \textbf{LDA} & \textbf{sLDA} \\
\midrule
Supervisione            & Non supervisionato & Supervisionato \\
Variabile risposta      & Nessuna       & $y_d$ (es.\ voto medio) \\
Obiettivo               & Massimizzare la verosimiglianza del corpus & Massimizzare congiuntamente corpus e $p(y_d \mid \theta_d)$ \\
Stima                   & Gibbs Sampling collassato & EM variazionale \\
Topic appresi per       & Coerenza testuale & Predittività di $y$ \\
Predizione out-of-sample & Non diretta & Sì, tramite $\hat{y} = \eta^\top \bar{\theta}$ \\
Interpretabilità        & Alta (topic tematici) & Media (topic più orientati alla risposta) \\
Complessità computazionale & Moderata   & Più elevata \\
\bottomrule
\end{tabular}
\end{table}

\paragraph{Limiti dell'sLDA.}
\begin{itemize}[leftmargin=1.8em]
  \item \textbf{Assunzione lineare}: il modello assume una
        relazione lineare tra $\bar{\theta}_d$ e $y_d$;
        relazioni non lineari non vengono catturate.
  \item \textbf{Numero di topic fisso}: come LDA, richiede la
        scelta di $K$ a priori.
  \item \textbf{Predizione su nuovi documenti}: richiede di
        rieseguire l'inferenza variazionale per ciascun nuovo
        documento; non esiste un metodo di predizione istantaneo
        come in un semplice modello lineare.
  \item \textbf{Overfitting in-sample}: le metriche
        (\emph{RMSE, MAE, $R^2$}) riportate dal pacchetto
        \pkg{lda} si riferiscono al fit sul training set;
        una corretta valutazione richiederebbe uno split
        train/test dedicato.
  \item \textbf{Sensibilità all'inizializzazione}: l'EM
        variazionale può convergere a ottimi locali; si
        raccomanda di eseguire più ripetizioni con semi diversi.
\end{itemize}

\subsection{Iperparametri e Risultati}

Il modello sLDA è stimato con gli stessi iperparametri dell'LDA
($K=6$, $\alpha=0.1$, $\eta=0.1$), con 100 iterazioni E-step
e 40 iterazioni M-step (\code{seed = 12345}).
La variabile risposta $y_d = \code{vote\_average}$ è trattata
come gaussiana (\code{logistic = FALSE}).

\begin{table}[H]
\centering
\caption{Topic sLDA: nomi assegnati e coefficienti di regressione}
\smallskip
\begin{tabular}{@{}lcc@{}}
\toprule
\textbf{Topic} & \textbf{Top-5 Parole} & \textbf{Coefficiente $\hat{\eta}$} \\
\midrule
Dramma Familiare \& Relazioni   & \ldots & \ldots \\
Crime \& Poliziesco             & \ldots & \ldots \\
Commedia Romantica \& Teen      & \ldots & \ldots \\
Horror \& Mystery               & \ldots & \ldots \\
Storico, Biografie \& Drama     & \ldots & \ldots \\
Fantascienza \& Fantasy         & \ldots & \ldots \\
\bottomrule
\end{tabular}
\end{table}

\noindent\textit{Nota: i valori esatti di top-words e
coefficienti dipendono dal seed e dai dati; sostituire con i
valori stampati da \code{top\_words\_slda} e
\code{slda\_model\$coefs} nell'output R.}

\subsection{Grafici sLDA}

\begin{figure}[H]
  \centering
  \includegraphics[width=\textwidth]{plot_slda_topwords.pdf}
  \caption{Le 10 parole con probabilità più alta per ciascun
           topic supervisionato (sLDA).}
  \label{fig:slda_topwords}
\end{figure}

\begin{figure}[H]
  \centering
  \begin{subfigure}[t]{0.48\textwidth}
    \includegraphics[width=\linewidth]{plot_slda_voto_topic.pdf}
    \caption{Distribuzione del voto per topic sLDA}
  \end{subfigure}
  \hfill
  \begin{subfigure}[t]{0.48\textwidth}
    \includegraphics[width=\linewidth]{plot_slda_pop_voto.pdf}
    \caption{Popolarità vs voto medio per topic sLDA}
  \end{subfigure}
  \caption{Relazione tra topic sLDA, gradimento e popolarità}
  \label{fig:slda_voto}
\end{figure}

\begin{figure}[H]
  \centering
  \includegraphics[width=0.80\textwidth]{plot_slda_coef.pdf}
  \caption{Coefficienti di regressione $\hat{\eta}$ per topic
           sLDA. Valori positivi indicano topic associati a
           voti più alti.}
  \label{fig:slda_coef}
\end{figure}

Il grafico dei coefficienti (Figura~\ref{fig:slda_coef})
permette di interpretare direttamente quali macro-temi narrativi
sono associati a voti più alti o più bassi: topic come
\textit{Storico \& Biografie} tendono ad avere coefficienti
positivi, mentre \textit{Horror \& Mystery} presentano
coefficienti più negativi, in accordo con quanto osservato nel
boxplot per genere.


% ================================================================
\section{Correlazione Parola-Voto}

Prima di costruire i modelli, si quantifica la correlazione
di Pearson tra la frequenza di ciascun termine nella DTM
e il voto medio \code{vote\_average}
(Figura~\ref{fig:correlazioni}).

\begin{figure}[H]
  \centering
  \includegraphics[width=0.85\textwidth]{plot_correlazioni.pdf}
  \caption{Top 15 termini con correlazione più positiva e
           più negativa rispetto a \code{vote\_average}}
  \label{fig:correlazioni}
\end{figure}

Le correlazioni sono in valore assoluto modeste
($|r| < 0.10$), anticipando i bassi valori di $R^2$ ottenuti
dai modelli: la sinossi da sola spiega una quota limitata
della varianza del voto.


% ================================================================
\section{Modelli di Regressione Penalizzata}

\subsection{Definizioni}

\begin{defbox}[Regressione Penalizzata]
I modelli di regressione penalizzata aggiungono alla funzione
di perdita ordinaria dei minimi quadrati un termine di
penalità sui coefficienti $\beta$:
\[
  \hat{\beta} = \arg\min_\beta \left[
    \sum_{i=1}^{n}(y_i - x_i^\top \beta)^2
    + \lambda \cdot P(\beta)
  \right]
\]
\begin{itemize}[leftmargin=1.8em, topsep=2pt]
  \item \textbf{Lasso} ($\alpha = 1$): $P(\beta) = \|\beta\|_1$,
        penalità $L_1$. Produce soluzioni \emph{sparse},
        annullando esattamente i coefficienti delle parole
        meno rilevanti --- de facto selezionando il vocabolario.
  \item \textbf{Ridge} ($\alpha = 0$): $P(\beta) = \|\beta\|_2^2$,
        penalità $L_2$. Mantiene tutti i predittori ma ne
        riduce l'ampiezza; utile con predittori altamente
        correlati.
  \item \textbf{Elastic Net} ($\alpha = 0.5$): combinazione
        lineare $\alpha\|\beta\|_1 + (1-\alpha)\|\beta\|_2^2$,
        che bilancia selezione e shrinkage.
\end{itemize}
Il parametro $\lambda$ è scelto tramite 10-fold
cross-validation come \code{lambda.min}.
\end{defbox}

\subsection{Costruzione della Document-Term Matrix}

La DTM è costruita con \pkg{tm} dalle sinossi preprocessate.
Si mantengono i termini presenti in almeno l'1\% dei documenti
(\code{removeSparseTerms(0.99)}), ottenendo una matrice di
\textbf{4\,867 film $\times$ 292~termini}.
I dati sono suddivisi in training~(75\%) e test~(25\%)
con \code{initial\_split}.

\subsection{Risultati}

\begin{table}[H]
\centering
\caption{Metriche sul test set -- Regressione penalizzata (DTM \pkg{tm})}
\smallskip
\begin{tabular}{@{}lccc@{}}
\toprule
\textbf{Modello} & \textbf{RMSE} & \textbf{MAE} & $\boldsymbol{R^2}$ \\
\midrule
Lasso       & 1.0100 & 0.8010 & 0.0575 \\
Ridge       & 1.0100 & 0.7990 & 0.0568 \\
Elastic Net & 1.0100 & 0.8010 & 0.0570 \\
\bottomrule
\end{tabular}
\end{table}

I tre modelli ottengono risultati pressoché identici con
$R^2 \approx 0.057$, indicando che le caratteristiche testuali
della sinossi spiegano circa il \textbf{5.7\%} della varianza
del voto.
Questo risultato, coerente con la letteratura sul topic,
riflette il fatto che il gradimento è determinato
principalmente da fattori non catturabili dalla sinossi
(qualità della regia, recitazione, marketing).

\begin{figure}[H]
  \centering
  \includegraphics[width=\textwidth]{plot_predicted_actual_penalizzata.pdf}
  \caption{Predicted vs Actual sul test set per i tre modelli
           penalizzati. La nuvola di punti attorno alla media
           riflette il basso potere predittivo della sinossi.}
  \label{fig:pred_act_pen}
\end{figure}

\begin{figure}[H]
  \centering
  \includegraphics[width=0.80\textwidth]{plot_coef_lasso_reg.pdf}
  \caption{Coefficienti non nulli del Lasso (regressione).
           Le parole in verde aumentano il voto previsto,
           quelle in rosso lo diminuiscono.}
  \label{fig:coef_lasso_reg}
\end{figure}


% ================================================================
\section{Modelli di Classificazione}

L'obiettivo è classificare automaticamente un film come
\textit{Alto} ($\geq 7$) o \textit{Basso} dalla sola sinossi.

\subsection{Sbilanciamento delle Classi}

Il dataset è fortemente sbilanciato:
\textbf{22.8\%} Alto vs \textbf{77.2\%} Basso
(1\,109 su 4\,858 film nella DTM sparsa).
Un classificatore banale (predice sempre \textit{Basso})
raggiungerebbe accuracy $\approx 77\%$ con sensitivity nulla.
Si adottano due strategie:

\begin{itemize}[leftmargin=1.8em]
  \item \textbf{LASSO Logistico}: cross-validation su AUC
        (anziché errore di classificazione), pesi per
        campione inversamente proporzionali alla frequenza di
        classe, soglia di classificazione fissata alla
        prevalenza nel training ($\approx 0.228$) anziché
        al default 0.5.
  \item \textbf{Random Forest}: pesi di classe
        \code{classwt = c(Alto=3, Basso=1)}.
\end{itemize}

\subsection{DTM Sparsa e Split}

La DTM sparsa (\code{cast\_sparse}, parole con
frequenza $\geq 20$) produce una matrice
\textbf{4\,858 film $\times$ 997~parole}.
Lo split train/test è stratificato 80\%/20\%
(\code{createDataPartition}).

\subsection{LASSO Logistico}

\begin{defbox}[LASSO Logistico]
Il LASSO Logistico è un modello di regressione logistica con
penalità $L_1$ applicata ai coefficienti. La funzione obiettivo
è:
\[
  \hat{\beta} = \arg\min_\beta \left[
    -\ell(\beta;\, X,\, y) + \lambda \|\beta\|_1
  \right]
\]
dove $\ell$ è la log-verosimiglianza binomiale. Come nel caso
lineare, la penalità $L_1$ induce sparsità, selezionando
automaticamente le parole più discriminanti tra le classi
\textit{Alto} e \textit{Basso}. La probabilità di appartenere
alla classe positiva è modellata tramite la funzione sigmoide:
$p(\text{Alto} \mid x) = \sigma(\beta^\top x)$.
\end{defbox}

Stimato con \code{cv.glmnet} (5-fold, $\alpha=1$,
famiglia \code{binomial}, misura AUC).
$\lambda_{\min} = 0.0026$.

\begin{figure}[H]
  \centering
  \begin{subfigure}[t]{0.48\textwidth}
    \includegraphics[width=\linewidth]{plot_roc_lasso.pdf}
    \caption{Curva ROC ($\mathrm{AUC} \approx 0.63$)}
  \end{subfigure}
  \hfill
  \begin{subfigure}[t]{0.48\textwidth}
    \includegraphics[width=\linewidth]{plot_coef_lasso_class.pdf}
    \caption{Coefficienti non nulli selezionati dal LASSO}
  \end{subfigure}
  \caption{LASSO Logistico: performance e parole discriminanti}
  \label{fig:lasso_class}
\end{figure}

Con la soglia adattiva il modello raggiunge una sensitivity
di \textbf{82.8\%} (individua correttamente la grande
maggioranza dei film di alta qualità), ma a prezzo di una
specificity molto bassa (\textbf{9.3\%}):
molti film a basso voto vengono erroneamente classificati
come \textit{Alto}, riflesso della natura ambigua del testo
delle sinossi.

\subsection{Random Forest}

\begin{defbox}[Random Forest]
Il \textbf{Random Forest} (Breiman, 2001) è un metodo di
\emph{ensemble learning} basato sulla combinazione di $B$
alberi decisionali addestrati su sottocampioni bootstrap del
dataset di training.

Ogni albero è costruito secondo il seguente meccanismo:
\begin{enumerate}[leftmargin=1.8em, topsep=2pt]
  \item Si estrae un campione bootstrap di $n$ osservazioni
        \emph{con rimpiazzo} dal training set.
  \item Ad ogni nodo dell'albero si seleziona casualmente un
        sottoinsieme di $m = \lfloor\sqrt{p}\rfloor$
        predittori (per classificazione) o
        $m = \lfloor p/3 \rfloor$ (per regressione), e si
        sceglie la variabile che massimizza il guadagno di
        informazione (Gini o varianza).
  \item Gli alberi crescono fino alla massima profondità
        (nessun pruning).
\end{enumerate}
La predizione finale è la \textbf{media} delle predizioni
dei singoli alberi (regressione) o il \textbf{voto a
maggioranza} (classificazione). La randomizzazione
a livello di campione e di feature riduce la varianza
rispetto al singolo albero (bias--variance tradeoff)
e rende il modello robusto all'overfitting.

L'\textbf{importanza} di una variabile è misurata come la
diminuzione media dell'accuratezza (\textit{Mean Decrease
Accuracy}) quando i valori di quella variabile vengono
permutati casualmente nelle osservazioni out-of-bag (OOB).
\end{defbox}

Addestrato sulla versione densa della DTM (top-300 parole,
300~alberi, $\mathrm{mtry} = \lfloor\sqrt{p}\rfloor$,
\code{classwt = c(Alto=3, Basso=1)}).

\begin{figure}[H]
  \centering
  \includegraphics[width=0.80\textwidth]{plot_rf_importance.pdf}
  \caption{Top 20 parole per importanza (Mean Decrease
           Accuracy) nel Random Forest -- classificazione}
  \label{fig:rf_imp}
\end{figure}

Il Random Forest raggiunge un \textbf{Kappa = 0.332},
indicando un accordo moderato al di là del caso.

\subsection{Confronto dei Classificatori}

\begin{table}[H]
\centering
\caption{Confronto dei modelli di classificazione sul test set}
\smallskip
\begin{tabular}{@{}lcccc@{}}
\toprule
\textbf{Modello} & \textbf{Accuracy} &
\textbf{Sensitivity} & \textbf{Specificity} & \textbf{Kappa} \\
\midrule
LASSO Logistico & 0.261 & 0.828 & 0.093 & $-$0.039 \\
Random Forest   & 0.733 & 0.606 & 0.770 & 0.332  \\
\bottomrule
\end{tabular}
\end{table}

I due modelli rivelano un chiaro \textbf{trade-off}: il LASSO
Logistico privilegia la sensitivity (rileva l'83\% dei film
di qualità), sacrificando la specificity; il Random Forest
offre un equilibrio migliore tra le due classi, come
confermato dal Kappa di Cohen più elevato.

% ================================================================
\section{Regressione sulla DTM Sparsa}

\begin{defbox}[LASSO Lineare sulla DTM Sparsa]
In questa sezione si utilizza la DTM sparsa (matrice
\textbf{4\,858 $\times$ 997}) anziché la DTM densa della
Sezione~\ref{sec:reg_pen}, includendo un vocabolario più ampio
(parole con frequenza $\geq 20$, contro la soglia dell'1\% dei
documenti). Il LASSO Lineare opera identicamente a quanto
descritto nella Sezione~\ref{sec:reg_pen}; la differenza risiede
nella maggiore ricchezza del vocabolario e nello split
stratificato 80/20.
\end{defbox}

Come confronto con la Sezione~\ref{sec:reg_pen}, si stima un
LASSO Lineare e un Random Forest sulla DTM sparsa
(stessa usata per la classificazione).

\begin{figure}[H]
  \centering
  \begin{subfigure}[t]{0.48\textwidth}
    \includegraphics[width=\linewidth]{plot_lasso_reg_sparse.pdf}
    \caption{LASSO Lineare}
  \end{subfigure}
  \hfill
  \begin{subfigure}[t]{0.48\textwidth}
    \includegraphics[width=\linewidth]{plot_rf_reg.pdf}
    \caption{Random Forest}
  \end{subfigure}
  \caption{Predicted vs Actual -- Regressione sulla DTM sparsa}
  \label{fig:reg_sparse}
\end{figure}

\begin{table}[H]
\centering
\caption{Confronto dei modelli di regressione (DTM sparsa,
         test set 20\%)}
\smallskip
\begin{tabular}{@{}lccc@{}}
\toprule
\textbf{Modello} & \textbf{RMSE} & \textbf{MAE} &
$\boldsymbol{R^2}$ \\
\midrule
LASSO Lineare  & 1.0531 & 0.8205 & 0.0819 \\
Random Forest  & 0.9554 & 0.7062 & 0.2444 \\
\bottomrule
\end{tabular}
\end{table}

Il Random Forest migliora sensibilmente rispetto al LASSO
Lineare ($R^2 = 0.244$ vs $0.082$, RMSE $= 0.955$ vs $1.053$),
suggerendo che la relazione tra il testo della sinossi e il
voto non è lineare e che le interazioni tra parole contengono
informazione predittiva non catturabile dalla regressione
penalizzata.


% ================================================================
\section{Discussione e Conclusioni}
\label{sec:conclusioni}

\subsection{Sintesi dei Risultati}

\begin{enumerate}[leftmargin=1.8em]
  \item \textbf{Vocabolario narrativo comune}: le sinossi
        condividono un nucleo lessicale ristretto; i bigrammi
        rivelano pattern narrativi (\textit{serial killer},
        \textit{high school}) inaccessibili al bag-of-words.

  \item \textbf{Topic LDA}: i 6~macro-temi latenti estratti
        corrispondono a generi riconoscibili e mostrano
        distribuzioni differenziate di voto e popolarità,
        evidenziando una tensione tra apprezzamento critico
        e successo commerciale.

  \item \textbf{sLDA}: l'estensione supervisionata orienta
        i topic verso la predizione del voto, producendo
        coefficienti $\hat{\eta}$ interpretabili che
        quantificano l'effetto di ciascun macro-tema sul
        gradimento del pubblico. L'$R^2$ in-sample è
        superiore ai modelli lineari, ma la valutazione
        out-of-sample rimane un limite metodologico da
        affrontare in sviluppi futuri.

  \item \textbf{Basso potere predittivo lineare}: tutti i
        modelli lineari convergono a $R^2 \approx 0.06$--$0.08$.
        Il Random Forest ottiene $R^2 = 0.244$, indicando
        che le non-linearità e le interazioni tra parole
        portano informazione aggiuntiva.

  \item \textbf{Classificazione e sbilanciamento}: il problema
        di classificazione è reso difficile dall'alto
        sbilanciamento (77\% vs 23\%). Il Random Forest con
        pesi di classe (Kappa $= 0.332$, sensitivity $= 61\%$,
        specificity $= 77\%$) offre il miglior compromesso;
        il LASSO Logistico con soglia adattiva massimizza la
        sensitivity (83\%) al costo di una specificity
        quasi nulla.
\end{enumerate}

\subsection{Limiti e Sviluppi Futuri}

\begin{itemize}[leftmargin=1.8em]
  \item Le sinossi sono prevalentemente in inglese; film
        in altre lingue hanno trame tradotte o assenti,
        introducendo un bias linguistico.
  \item Il modello usa esclusivamente la sinossi. Integrare
        variabili strutturali (\code{vote\_count},
        \code{year}, \code{genre}) aumenterebbe
        sensibilmente il potere predittivo.
  \item Rappresentazioni semantiche dense
        (\textit{word embeddings}, Sentence-BERT) potrebbero
        catturare la semantica in modo più ricco rispetto
        al bag-of-words.
  \item Per la classificazione, tecniche di oversampling
        (SMOTE) o modelli di ensemble specializzati per
        classi sbilanciate potrebbero migliorare
        ulteriormente il Kappa.
  \item Per l'sLDA, un corretto split train/test permetterebbe
        di valutare il potere predittivo out-of-sample e
        confrontarlo direttamente con gli altri modelli.
\end{itemize}

% ================================================================
\end{document}
